\section{Finite morphisms}
\label{section-finite-morphisms}

\noindent
In this section we collect some observations about the
constructions in the previous sections for finite morphisms.
Let $f : Y \to X$ be a finite morphism of locally Noetherian schemes.
Let $\omega_{Y/X}$ be as in
Remark \ref{remark-relative-dualizing-for-quasi-finite}.

\medskip\noindent
The first remark is that
$$
f_*\omega_{Y/X} = \SheafHom_{\mathcal{O}_X}(f_*\mathcal{O}_Y, \mathcal{O}_X)
$$
as sheaves of $f_*\mathcal{O}_Y$-modules. Since $f$ is affine, this
formula uniquely characterizes $\omega_{Y/X}$, see
Morphisms, Lemma \ref{morphisms-lemma-affine-equivalence-modules}.
The formula holds because for $\Spec(A) = U \subset X$ affine open, the
inverse image $V = f^{-1}(U)$ is the spectrum of a finite $A$-algebra
$B$ and hence
$$
H^0(U, f_*\omega_{Y/X}) =
H^0(V, \omega_{Y/X}) =
\omega_{B/A} =
\Hom_A(B, A) =
H^0(U, \SheafHom_{\mathcal{O}_X}(f_*\mathcal{O}_Y, \mathcal{O}_X))
$$
by construction. In particular, we obtain a canonical evaluation map
$$
f_*\omega_{Y/X} \longrightarrow \mathcal{O}_X
$$
which is given by evaluation at $1$ if we think of $f_*\omega_{Y/X}$
as the sheaf $\SheafHom_{\mathcal{O}_X}(f_*\mathcal{O}_Y, \mathcal{O}_X)$.

\medskip\noindent
The second remark is that using the evaluation map we obtain
canonical identifications
$$
\Hom_Y(\mathcal{F}, f^*\mathcal{G} \otimes_{\mathcal{O}_Y} \omega_{Y/X})
=
\Hom_X(f_*\mathcal{F}, \mathcal{G})
$$
functorially in the quasi-coherent module $\mathcal{F}$ on $Y$
and the finite locally free module $\mathcal{G}$ on $X$.
If $\mathcal{G} = \mathcal{O}_X$ this follows immediately
from the above and
Algebra, Lemma \ref{algebra-lemma-adjoint-hom-restrict}.
For general $\mathcal{G}$ we can use the same lemma and the
isomorphisms
$$
f_*(f^*\mathcal{G} \otimes_{\mathcal{O}_Y} \omega_{Y/X}) =
\mathcal{G} \otimes_{\mathcal{O}_X}
\SheafHom_{\mathcal{O}_X}(f_*\mathcal{O}_Y, \mathcal{O}_X) =
\SheafHom_{\mathcal{O}_X}(f_*\mathcal{O}_Y, \mathcal{G})
$$
of $f_*\mathcal{O}_Y$-modules where the first equality is the
projection formula
(Cohomology, Lemma \ref{cohomology-lemma-projection-formula}).
An alternative is to prove the formula affine locally by
direct computation.

\medskip\noindent
The third remark is that if $f$ is in addition flat, then the
composition
$$
f_*\mathcal{O}_Y \xrightarrow{f_*\tau_{Y/X}} f_*\omega_{Y/X}
\longrightarrow \mathcal{O}_X
$$
is equal to the trace map $\text{Trace}_f$ discussed in
Section \ref{section-discriminant}. This follows immediately by
looking over affine opens.

\medskip\noindent
The fourth remark is that if $f$ is flat and $X$ Noetherian, then
we obtain
$$
\Hom_Y(K, Lf^*M \otimes_{\mathcal{O}_Y} \omega_{Y/X})
=
\Hom_X(Rf_*K, M)
$$
for any $K$ in $D_\QCoh(\mathcal{O}_Y)$ and $M$ in $D_\QCoh(\mathcal{O}_X)$.
This follows from the material in
Duality for Schemes, Section \ref{duality-section-proper-flat},
but can be proven directly in this case as follows.
First, if $X$ is affine, then it holds by
Dualizing Complexes, Lemmas \ref{dualizing-lemma-right-adjoint} and
\ref{dualizing-lemma-RHom-is-tensor-special}\footnote{There is a
simpler proof of this lemma in our case.} and
Derived Categories of Schemes, Lemma \ref{perfect-lemma-affine-compare-bounded}.
Then we can use the induction principle
(Cohomology of Schemes, Lemma \ref{coherent-lemma-induction-principle})
and Mayer-Vietoris
(in the form of Cohomology, Lemma \ref{cohomology-lemma-mayer-vietoris-hom})
to finish the proof.







